% STRATA condensed paper: built section by section from paper-v2.tex during the 2026-09-29 paragraph-by-paragraph review.
% Each section here is the condensed form of the corresponding section of the long manuscript; numbers are identical.
\documentclass[11pt,letterpaper]{article}

\usepackage[margin=1in]{geometry}
\usepackage{microtype}
\usepackage[hidelinks]{hyperref}
\usepackage{cite}
\usepackage{times}
\usepackage{graphicx}
\usepackage{array}
\usepackage{longtable}

\title{STRATA: A Stratified Process-Mining Framework for HVAC Fault Detection and Diagnosis}
\author{Yassh Singh \\
\small Supervisors: Dr.\ Anthony Aighobahi, Dr.\ Mridula Sharma \\
\small Thompson Rivers University --- UREAP 2026}
\date{September 2026}

\begin{document}
\maketitle

\begin{abstract}
\noindent
Faults in building HVAC equipment persist because the tools meant to catch them either raise alarms that operators learn to ignore, or flag anomalies that nobody can act on. STRATA turns ordinary one-minute trend data into an event log along the control sequence, discovers process models at the device and unit level, and scores each day through eight calibrated, physically interpretable channels against a measured noise floor, so that every alarm names a location, states its evidence and is held to a false-alarm budget. It was built under a falsification protocol: every hypothesis was pre-registered with the observation that would refute it, every result was recomputed independently before it was quoted, and automated tests pin the published numbers and the falsifiers that fired.

On five public LBNL benchmark systems, two of them never previously benchmarked, STRATA detects 158 of 175 seeded faults, most on the first day, at 26 false-alarm days in 480 held-out fault-free days (1 to 9.4\% per system). Two systems were onboarded blind, by configuration alone. The alarm names the correct terminal unit in 37 of the 50 detections with known ground truth and a wrong one in none, and the faulted subsystem in 83\% of all detections. The false-alarm budget holds under injected sensor error, change-of-value logging and schedule shifts (worst 7 of 96 days) and on one measured rooftop unit (3 of 49).

Process mining helps to get detection but is not detecting on its own. The strata and the event abstraction carry the localisation, and invariants that discovery surfaced became several of the rules that detect. Conformance checking of the discovered models, pre-registered as the detector, was refuted: over 199 scenarios on six systems, no fault is caught by the models alone. A principal-component baseline detects three fewer of the 73 scenarios it shares and cannot name the device. A third-party Guideline 36 rule battery, held to the same gate, matches STRATA within one detection on the air handler and detects a third as many faults on the terminal units, so the pre-registered improvement claim over current practice is withdrawn there. The evaluation also produced eight machine-verified errata in the benchmark datasets. Every number regenerates from the public repository.
\end{abstract}

% Sections follow as the review proceeds.

\bibliographystyle{plain}
\bibliography{../references}
\end{document}
